% ============================================================
%  Chapter 11 – Further Techniques and Future Directions
% ============================================================

\section{Further Techniques and Future Directions}%
\label{sec:future-directions}

Descriptions of external algorithms in this chapter are based on cited
sources and may not reflect the latest versions of those tools.

The deep-zoom Mandelbrot rendering community has developed a rich
collection of techniques that complement or extend the algorithms
implemented in \FractalShark{}.  Understanding these methods is valuable
for two reasons: several of them directly influenced \FractalShark{}'s
design choices (e.g.\ the decision to use rebasing rather than
glitch correction, or linear approximation rather than series
approximation), and others represent promising avenues for future
improvement.

Each subsection below describes a technique, gives its key formulas, and
explains how it relates to---and differs from---\FractalShark{}'s current
approach.  \Cref{fig:design-decisions} summarizes the three principal
design choices at a glance.

\begin{figure}[!htbp]
\centering
\resizebox{\linewidth}{!}{%
\begin{tikzpicture}[
  every node/.style={font=\small},
  problem/.style={draw, rounded corners, fill=blue!8, minimum width=3.8cm,
                  minimum height=0.7cm, align=center},
  option/.style={draw, rounded corners, minimum width=3.6cm,
                 minimum height=0.6cm, align=center},
  chosen/.style={option, fill=orange!15, line width=1pt},
  why/.style={font=\scriptsize\itshape, text width=3.2cm, align=center},
  arr/.style={-{Stealth[length=5pt]}, thick}
]

% --- Column 1: Numerical stability ---
\node[problem] (p1) {Orbit approximation\\errors};
\node[option,  below left =0.8cm and -0.1cm of p1]  (a1) {Glitch detection\\+ correction};
\node[chosen,  below right=0.8cm and -0.1cm of p1]  (b1) {Rebasing};
\node[why, below=0.15cm of a1] {multi-pass,\\post-hoc};
\node[why, below=0.15cm of b1] {single-pass,\\inline};

\draw[arr] (p1.south west) ++(0.3,0) -- (a1.north);
\draw[arr] (p1.south east) ++(-0.3,0) -- (b1.north);

% --- Column 2: Iteration skipping ---
\node[problem, right=2.8cm of p1] (p2) {High iteration\\count};
\node[option,  below left =0.8cm and -0.1cm of p2]  (a2) {Series Approx.\\(SA)};
\node[chosen,  below right=0.8cm and -0.1cm of p2]  (b2) {Linear Approx.\\(LA~v2)};
\node[why, below=0.15cm of a2] {Taylor series,\\divergence risk};
\node[why, below=0.15cm of b2] {adaptive spans,\\validity thresholds};

\draw[arr] (p2.south west) ++(0.3,0) -- (a2.north);
\draw[arr] (p2.south east) ++(-0.3,0) -- (b2.north);

% --- Column 3: Multiplication ---
\node[problem, right=2.8cm of p2] (p3) {$O(n^2)$\\multiplication};
\node[option,  below left =0.8cm and -0.1cm of p3]  (a3) {Karatsuba /\\Toom--Cook};
\node[chosen,  below right=0.8cm and -0.1cm of p3]  (b3) {NTT};
\node[why, below=0.15cm of a3] {recursive,\\$O(n^{1.58})$};
\node[why, below=0.15cm of b3] {GPU-parallel,\\$O(n\log n)$};

\draw[arr] (p3.south west) ++(0.3,0) -- (a3.north);
\draw[arr] (p3.south east) ++(-0.3,0) -- (b3.north);

\end{tikzpicture}%
}% end resizebox
\caption[Algorithm design decisions]{Three principal design decisions
  in \FractalShark{}.  Shaded boxes indicate the chosen approach;
  italic annotations note the key trade-off in each case.}
\label{fig:design-decisions}
\end{figure}

% -----------------------------------------------------------
\subsection{Glitch Detection and Correction}
\label{subsec:glitch-detection}

Finite-precision perturbation can produce \emph{glitches}: incorrect pixels
caused by an unsuitable reference or loss of information in the delta
representation.  A cancellation test can identify stressed pixels for
another reference calculation.  One remedy is a multi-pass
detect-and-correct workflow.

The \emph{Pauldelbrot criterion} flags possible glitches when cancellation
makes the reconstructed pixel state small relative to the reference:

\begin{equation}
  |z_n^\star+\Delta z_n|^2 < G\,|z_n^\star|^2,\qquad 0<G<1,
  \label{eq:pauldelbrot-criterion}
\end{equation}
where $G$ is an empirical sensitivity setting.  For a real reference
$z_n^\star=1$ and delta $\Delta z_n=-0.999$, the reconstructed state is
$0.001$.  With $G=10^{-4}$, its squared magnitude $10^{-6}$ is below
$10^{-4}$, so the test fires even though $|\Delta z_n|<|z_n^\star|$.
Comparing reference and delta magnitudes alone misses the cancellation.
The criterion is a diagnostic heuristic, not a proof that every flagged
pixel is wrong.  The correction procedure is to scan the rendered image for such pixels,
compute a new reference orbit centered on a glitched location, and
re-render only the affected region.  This process repeats until no
glitches remain and requires multi-reference orbit support
(\cref{subsec:multi-reference}).

\FractalShark{} uses \emph{rebasing} (\cref{sec:perturb-only-rebase})
within pixel iteration.  When $|\Delta z|>|z|$, where $z$ is the
reconstructed current pixel state, it sets $\Delta z\leftarrow z$ and
resets the reference index to zero while retaining the accumulated pixel
iteration count.  In the example above, the replacement delta is $0.001$.
The algebra preserves $z$ because the new reference value is $z_0^\star=0$.
Rebasing reduces this representation's cancellation problem, but cannot
recover information already rounded away, repair an inaccurate reference, or
certify LA/BLA truncation error.  Glitch diagnostics and additional references
can therefore remain useful validation or fallback strategies.
The cancellation criterion and rebasing are discussed in
\cite{HeilandAllen2021DeepZoom,HeilandAllen2022DeepZoom}.

% -----------------------------------------------------------
\subsection{Series Approximation (SA)}
\label{subsec:series-approx}

Series Approximation was the original technique for skipping iterations
in perturbation-based Mandelbrot rendering, introduced alongside
perturbation theory itself~\cite{Martin2013Perturbation}.  \FractalShark{} instead uses BLA~v1
(\cref{sec:bilinear-approx}) and LA~v2
(\cref{sec:la-v2-perturb}).  SA illustrates another way to trade
precomputation for per-pixel work; performance depends on the view and the
chosen approximation order.

SA constructs a polynomial $P(\Delta c)$ that approximates the
perturbation $\Delta z_N$ after $N$~iterations, so that each pixel can
skip those iterations with a single polynomial evaluation.  Starting
from the perturbation recurrence
\begin{equation}
  \Delta z_{n+1} = 2\,z_n^\star\,\Delta z_n + (\Delta z_n)^2 + \Delta c,
  \label{eq:sa-perturb-recurrence}
\end{equation}
one expands $\Delta z_n$ as a power series in $\Delta c$:
\[
  \Delta z_n = A_n\,\Delta c + B_n\,(\Delta c)^2
             + C_n\,(\Delta c)^3 + \cdots
\]
Substituting into \cref{eq:sa-perturb-recurrence} and matching powers of
$\Delta c$ yields recurrences for the coefficients:
\begin{align}
  A_{n+1} &= 2\,z_n^\star\,A_n + 1,
    \label{eq:sa-A} \\
  B_{n+1} &= 2\,z_n^\star\,B_n + A_n^2,
    \label{eq:sa-B} \\
  C_{n+1} &= 2\,z_n^\star\,C_n + 2\,A_n\,B_n,
    \label{eq:sa-C}
\end{align}
with $A_0=B_0=C_0=0$ and higher-order terms following the same pattern.
The linear coefficient $A_n$ is the derivative $\partial z_n/\partial c$.
The parameter coefficient in a composed LA map captures the corresponding
first-order parameter dependence, although LA also retains a nonlinear input
operation (\cref{sec:la-step-forward}).

For $c_\star=-1$, the first coefficients are
$(A_1,B_1,C_1)=(1,0,0)$,
$(A_2,B_2,C_2)=(-1,1,0)$, and
$(A_3,B_3,C_3)=(1,1,-2)$.
Thus a cubic approximation after three steps is
$\Delta z_3\approx\Delta c+(\Delta c)^2-2(\Delta c)^3$.
The exact expression has one additional term $(\Delta c)^4$.
At $\Delta c=0.001$, these give $0.001000998$ and
$0.001000998001$, respectively, showing both how skipping works and what
truncation discards.

\paragraph{Why the implementation uses linear approximation.}
SA and linear approximation (BLA~v1, LA~v2) address the same
problem---skipping iterations that would otherwise be computed one at a
time---but they differ fundamentally in approach.  SA uses a
higher-order polynomial in $\Delta c$: a single evaluation can skip many
iterations, but a truncated polynomial needs a remainder estimate or other
applicability test.  The formulas above alone provide neither.  Implementations
must select a suitable order and skip length and fall back when the retained
polynomial is insufficient.

Linear approximation, by contrast, uses first-order (linear)
coefficients with explicit, precomputed validity thresholds
(\cref{sec:la-threshold}). BLA~v1 (\cref{sec:bilinear-approx}) builds a
generic power-of-two hierarchy of composable linear maps, agnostic of
the orbit's dynamics. LA~v2 (\cref{sec:la-v2-perturb}) retains a
nonlinear input operation and chooses adaptive boundaries near useful
returns, with synthetic spans when no suitable return is detected.
Higher stages compose consecutive entries; neither stage-0 spans nor
their compositions must be multiples of a fundamental period.
BLA~v1 falls back to individual perturbation updates when no retained
map is usable. LA~v2 descends through finer stages before handing
remaining work to its perturbation tail (\cref{sec:la-v2-motivation}).

The linear approach avoids higher-order coefficient management and
checks block applicability before evaluation. Useful return spans can
enable very large skips near periodic points, but the thresholds do
not certify a complete orbit error bound or a universal performance
advantage. Optimized SA variants such as NanoMB1,
NanoMB2, and Super Series Approximation (SSA)
(\cite{KallesFraktaler}) use more elaborate coefficient and validity
management.  These alternatives do not establish a universal speed ordering.

% -----------------------------------------------------------
\subsection{Multi-Reference Orbit Strategies}
\label{subsec:multi-reference}

In renderers that use the glitch-detection workflow described in
\cref{subsec:glitch-detection}, a single reference orbit may not serve
all pixels well.  The multi-reference approach addresses this by
computing reference orbits at multiple locations and assigning each pixel
to the best available reference---the one that minimizes $|\Delta z_n|$
at the iteration where the perturbation is most stressed.

Multi-reference rendering typically proceeds iteratively:
\begin{enumerate}
  \item Render the image with the primary reference orbit.
  \item Detect glitched pixels (\cref{subsec:glitch-detection}).
  \item Compute a new reference orbit centered on a glitched pixel's
        location.
  \item Re-render only the glitched pixels using the new reference.
  \item Repeat until no glitches remain.
\end{enumerate}

Per-pixel rebasing (\cref{sec:perturb-only-rebase}) reduces the need for
extra references in \FractalShark{}, but does not remove every possible
benefit.  A reference nearer a difficult region can reduce delta magnitudes,
improve approximation applicability, or help investigate numerical disagreement.
The additional reference computation, memory, and pixel assignment must be
weighed against that benefit.  The approach is discussed in
\cite{HeilandAllen2022DeepZoom}.

% -----------------------------------------------------------
\subsection{Smooth (Continuous) Iteration Count}
\label{subsec:smooth-coloring}

Standard escape-time coloring uses integer iteration counts, producing
visible banding in the rendered image.  The \emph{smooth iteration
count} eliminates this artifact by computing a fractional escape count.

One common convention for the smooth count is
\begin{equation}
  n_{\text{smooth}} = n + 1
    - \frac{\log\!\bigl(\log|z_n|\bigr)}{\log 2},
  \label{eq:smooth-iter}
\end{equation}
where $n$ counts completed recurrence steps and $z_n$ is the corresponding
escaped orbit value.  The $+1$ and any other additive constant define the
palette's indexing convention; they do not affect the removal of bands.
The derivation rests on the potential function
$G(c) = \lim_{k\to\infty} \log|z_k| / 2^k$, which is continuous across
the exterior of the Mandelbrot set.  The escape time is related to the
level set $G(c) = \log R / 2^n$; interpolating between successive
iterates yields the fractional part.

A threshold-aligned convention for a general bailout radius $R>1$ is
\begin{equation}
  n_{\text{smooth},R} = n
    - \frac{\log\!\bigl(\log|z_n|/\log R\bigr)}{\log 2}.
  \label{eq:smooth-iter-general}
\end{equation}

The expression gives $n$ when $|z_n|=R$ and differs from
\eqref{eq:smooth-iter} by a constant.  To derive it, neglect $c$ after
escape so that $\log|z_{n+k}|\approx2^k\log|z_n|$, then solve for the
fractional time at which this growth would cross $R$.
For $R=2$ and $|z_n|=4$, it gives $n-1$.  One further idealized square
produces $|z_{n+1}|=16$, giving $(n+1)-2=n-1$ again: the smoothed
count is approximately independent of which later iterate is sampled.
Using a larger bailout makes the neglected $c$ term less important.

Fractional counts can drive interpolated palette lookup
(\cref{subsec:coloring}) and supersampled coloring
(\cref{subsec:antialiasing}).  The current render output contains iteration
counts, not each pixel's escaped $|z_n|^2$.  Adding smooth coloring therefore
requires computing the fractional value in the iteration kernel or extending
its output to retain the escaped magnitude; it cannot be recovered from the
existing count alone.  Smooth coloring is discussed in
\cite{HeilandAllenHybridDistance,WikipediaPlottingAlgorithms}.

% -----------------------------------------------------------
\subsection{Distance Estimation}
\label{subsec:distance-estimation}

Exterior distance estimation relates orbit growth to parameter sensitivity,
providing a useful scale for deciding whether a pixel is near the boundary.
A commonly used potential-over-gradient distance scale is
\begin{equation}
  d(c) \;\approx\; \frac{|z_n|\,\log|z_n|}{|z'_n|},
  \label{eq:distance-estimate}
\end{equation}
where $z'_n = \partial z_n / \partial c$ is the \emph{pixel} orbit
derivative after escape.  Taking $G\approx2^{-n}\log|z_n|$ gives a
gradient magnitude approximately $2^{-n}|z'_n|/|z_n|$; dividing $G$ by
that magnitude cancels $2^{-n}$ and yields \eqref{eq:distance-estimate}.
The scale is not the exact shortest distance.  Other conventions multiply it
by two to approximate an upper distance bound; rigorous bounds require their
stated analytic conditions and a sufficiently accurate escaped orbit
\cite{HeilandAllenExteriorDistance}.
For $c=1$, $z_3=5$ and $z'_3=13$, giving a scale
$5\log5/13\approx0.619$ in parameter-plane units.

\FractalShark{} already implements derivative recurrences for reference
periodicity (\cref{subsec:ref-orbit-periodicity}) and Feature Finder
approximation evaluation (\cref{sec:la-derivatives}).  Those are not escaped
per-pixel derivatives in the normal render output.  Distance rendering would
need to propagate a derivative along each pixel's accepted approximation and
perturbation steps, preserve its scaling through rebases, and consume it with
the escaped magnitude.  The additional arithmetic and output state are part
of the feature, along with coloring integration.

Applications of the distance estimator include:
\begin{itemize}
  \item Anti-aliased boundary rendering, where pixels near $\partial
        \mathcal{M}$ receive a smooth opacity falloff.
  \item Resolution-independent thin-filament rendering, ensuring that
        fine structures are visible regardless of pixel density.
  \item 3D height-field visualization, using $d(c)$ as a height value
        to produce surface renders of the set.
\end{itemize}

Interior distance estimation is also possible—using the Koebe quarter
theorem and conformal mapping to the unit disk—but is substantially more
complex and requires detecting the attracting cycle first.  The
interaction of perturbation rendering with distance estimation is
discussed in \cite{HeilandAllen2013Perturbation}.

% -----------------------------------------------------------
\subsection{Interior Rendering}
\label{subsec:interior-rendering}

Pixels that do not escape within the selected iteration budget are
traditionally rendered as uniform black.  That finite observation does not
prove membership in the Mandelbrot set.  For parameters whose critical orbit
converges to an attracting cycle, the cycle's period and multiplier provide
additional interior structure.  Such parameters belong to hyperbolic
components; arbitrary bounded orbits, especially boundary orbits, cannot all
be described that way.

Several interior coloring techniques exploit this structure:
\begin{itemize}
  \item \emph{Period-based coloring.}  Detect the period~$p$ of the
        attracting cycle and assign a color according to~$p$.  This
        reveals the intricate structure of hyperbolic components.
  \item \emph{Phase coloring.}  Use the argument
        $\arg(z^\star)$ of the attracting fixed point~$z^\star$ to
        produce smooth color gradients within a single component.
  \item \emph{Interior distance.}  Estimate how far the point lies from
        the boundary of its component, analogous to the exterior
        distance estimator (\cref{subsec:distance-estimation}).
\end{itemize}

For a cycle $z_0^*,\ldots,z_{p-1}^*$, the multiplier is
$\lambda=\prod_{k=0}^{p-1}2z_k^*$.  A small displacement after one
cycle is approximately multiplied by $\lambda$, so $|\lambda|<1$
explains attraction.  At $c=0$, the fixed point $0$ has multiplier zero;
at $c=-1$, the cycle $0,-1$ also has multiplier zero.  Both remain
unescaped, but period coloring distinguishes them.

Floyd's or Brent's algorithms detect exact repetition in a discrete state
sequence with $O(1)$ extra memory.  A finite-precision orbit can repeat
without proving a mathematical cycle; numerical attracting-cycle detection
needs tolerance and multiplier checks.  Near a neutral cycle, convergence
can be too slow for a short iteration budget.

\FractalShark{}'s Feature Finder (\cref{sec:feature-finder}) already
detects periods via Newton's method on the equation $z_p(c) = 0$; the
same period information could be used for interior coloring if the
detection logic were extended to run during rendering rather than as a
separate offline tool.

% -----------------------------------------------------------
\subsection{Boundary Tracing}
\label{subsec:boundary-tracing}

In a typical Mandelbrot image, large contiguous regions share the same
integer iteration count.  \emph{Boundary tracing} exploits this
observation: instead of computing every pixel independently, the
algorithm traces the contours between regions of different iteration
counts and flood-fills the interiors.

The procedure for each connected region of constant iteration count is:
\begin{enumerate}
  \item Identify a border pixel (one whose neighbor has a different
        iteration count).
  \item Follow the 4- or 8-connected border, recording the contour.
  \item Fill all enclosed pixels with the shared iteration count,
        skipping the per-pixel iteration entirely for interior pixels.
\end{enumerate}

The speedup depends on the fraction of the image occupied by large uniform
regions and on how much verification is needed.  A uniform outer border does
not prove a uniform interior: a $3\times3$ grid with every border count equal
to $5$ can still have center count $20$.  Mandelbrot images can contain
similarly enclosed filaments or islands.  A reliable implementation therefore
needs additional interior probes, subdivision, or region-specific bounds
before filling unsampled pixels.

\paragraph{Trade-offs.}
Flood-filling one constant count suits integer escape-time coloring.
For smooth coloring (\cref{subsec:smooth-coloring}), equal integer counts
can still have different fractional values, so the same fill would erase
those gradients.  Adaptive region methods can be extended to continuous
values, but require an interpolation error test instead of constant-count
agreement.

Boundary tracing is orthogonal to perturbation and LA: it is a
pixel-level optimization applied \emph{after} per-pixel iteration counts
are available for the border pixels; the per-pixel algorithm
(perturbation, LA, etc.)\ is unchanged.  Sampling and fill validation remain
additional responsibilities.  An overview of the algorithm is given in
\cite{WikipediaPlottingAlgorithms}.

% -----------------------------------------------------------
\subsection{Histogram Equalization}
\label{subsec:histogram-equalization}

At deep zooms, iteration count distributions are highly skewed: most
pixels cluster around similar counts, with a long tail of rare high
values.  Standard palette cycling maps iteration counts to colors via
modular indexing, which wastes most of the color palette on the
infrequent tail values and compresses the visually dominant range into a
narrow band.

\emph{Histogram equalization} addresses this by redistributing the
palette according to the actual distribution of iteration counts:
\begin{enumerate}
  \item Collect a histogram $H(k)$ of iteration counts across all
        pixels.
  \item Compute the cumulative distribution function
        $\mathrm{CDF}(k) = \sum_{j \le k} H(j) \big/ \sum_j H(j)$.
  \item Map each pixel's iteration count~$k$ to the color at palette
        index $\mathrm{CDF}(k)$.
\end{enumerate}
For example, four exterior pixels with count $10$, two with count $20$, and
two with count $50$ give $\mathrm{CDF}(10)=1/2$,
$\mathrm{CDF}(20)=3/4$, and $\mathrm{CDF}(50)=1$.
All four pixels tied at $10$ retain the same mapped value: equalization cannot
split one repeated count across a range of colors without another rule.
It spreads distinct values according to their observed frequencies rather
than guaranteeing uniform palette occupancy.  Unescaped pixels normally
receive a separate interior color and are excluded from this exterior CDF.

\FractalShark{} already computes per-frame iteration statistics
(minimum, maximum, sum) via the reduction kernel
(\cref{subsec:reduction}); extending this to a full histogram would
require histogram collection and a second coloring pass.  Wide iteration
ranges may require binned or sparse histograms rather than one counter per
possible count.  A changing per-frame CDF can also change colors during a
zoom animation, motivating a stable shared distribution.  The technique is
independent of the iteration algorithm and can be combined
with smooth coloring (\cref{subsec:smooth-coloring}) for even better
results.  A general discussion of histogram-based coloring appears in
\cite{WikipediaPlottingAlgorithms}.

The techniques surveyed in this chapter represent a broader landscape of
fractal rendering optimizations beyond those currently implemented in
\FractalShark{}.  Of the approaches discussed, rebasing (over glitch
correction), linear approximation (over series approximation), and
NTT-based multiplication (over Karatsuba) have directly shaped the
system's present architecture, as detailed in the preceding chapters.
The remaining techniques---smooth coloring, distance estimation,
interior rendering, boundary tracing, and histogram
equalization---represent possible extensions with different state,
verification, and synchronization costs.  Smooth and distance coloring need
additional per-pixel quantities, whereas boundary tracing must justify filling
unsampled regions.
\Cref{sec:conclusion} summarizes the overall design and reflects on the
decisions that connect these subsystems into a coherent whole.
